\documentclass[../../../include/open-logic-section]{subfiles}

\begin{document}

\olfileid[hi]{sth}{story}{extensionality}
\olsection{विस्तारिता}

सबसे पहले यह कहना चाहिए कि समुच्चयों की पहचान उनके !!{element}s से होती है।
अधिक सटीक रूप में:

\begin{axiom}[विस्तारिता]
यदि समुच्चयों $A$ और $B$ के !!{element}s समान हैं, तो $A$ और $B$ एक ही
समुच्चय हैं।
\[
  \lforall[A][\lforall[B][(\lforall[x][(x \in A \liff x \in B)] \lif
  \eq[A][B])]]
\]
\end{axiom}

हमने पूरे \olref[sfr][][]{part} में इसे मान लिया था, किंतु इसे दोहराना
उपयोगी है। विस्तारिता का अभिगृहीत इस मूल विचार को व्यक्त करता है कि किसी
समुच्चय का निर्धारण उसके !!{element}s से होता है। (इसलिए समुच्चयों की तुलना
\emph{संकल्पनाओं} से की जा सकती है, जहाँ ठीक वही वस्तुएँ अनेक भिन्न
संकल्पनाओं के अंतर्गत आ सकती हैं।)

इस सिद्धांत को क्यों स्वीकार करें? यह कहना युक्तिसंगत है कि विस्तारिता का
कोई भी निषेध उस हर चीज़ को त्यागने का निर्णय है जिसे \emph{समुच्चय सिद्धांत}
तक कहा जा सके। समुच्चय सिद्धांत विस्तारात्मक संग्रहों का सिद्धांत ही है---न
उससे अधिक, न कम।

फिर भी \olref[sth][][]{part} की वास्तविक चुनौती ऐसे सिद्धांत निर्धारित
करना है जो हमें बताएँ कि \emph{कौन-से समुच्चय अस्तित्व रखते हैं}। और पता
चलता है कि इस प्रश्न का एकमात्र सचमुच ``स्पष्ट'' उत्तर सिद्धतः गलत है।

\end{document}
